\begin{theorembox}{Theorem (Uniqueness of Limits)}
% ---------------------------------------------------------
% Theorem: Uniqueness of Limits
% Label: thm:uniqueness-of-limits
% ---------------------------------------------------------
\begin{theorem}[Uniqueness of Limits]
\label{thm:uniqueness-of-limits}

Let $(x_n) \subseteq \mathbb{R}$. If $x_n \to L$ and $x_n \to K$, then
$L = K$.

\smallskip
\noindent
\hyperref[prf:uniqueness-of-limits]{\textit{Go to proof.}}
\end{theorem}
\end{theorembox}

\begin{remark*}[Standard quantified statement]
\[
\begin{aligned}
&\forall (x_n) \subseteq \mathbb{R} \;\forall L \in \mathbb{R}
\;\forall K \in \mathbb{R}\\
&\quad\Bigl(
    \bigl(x_n \to L \;\land\; x_n \to K\bigr)
    \;\Longrightarrow\;
    L = K
\Bigr).
\end{aligned}
\]
\end{remark*}

\begin{remark*}[Predicate reading]
\[
\bigl(\operatorname{ConvergesTo}(x_n,L,\mathbb{R})
\land
\operatorname{ConvergesTo}(x_n,K,\mathbb{R})\bigr)
\Longrightarrow L=K.
\]
\end{remark*}

\begin{remark*}[Negated quantified statement]
\[
\begin{aligned}
&\exists (x_n) \subseteq \mathbb{R} \;\exists L \in \mathbb{R}
\;\exists K \in \mathbb{R}\\
&\quad\Bigl(
    \bigl(x_n \to L \;\land\; x_n \to K\bigr)
    \;\land\;
    L \neq K
\Bigr).
\end{aligned}
\]
\end{remark*}

\begin{remark*}[Negation predicate reading]
\[
\operatorname{ConvergesTo}(x_n,L,\mathbb{R})
\land
\operatorname{ConvergesTo}(x_n,K,\mathbb{R})
\land L\neq K.
\]
\end{remark*}
\begin{remark*}[Failure modes]
\begin{description}
  \item[Exposition.]
Failure is witnessed by the negated condition displayed above.
  \item[\operatorname{ConvergesTo}.]
\[
\begin{aligned}
&\exists (x_n) \subseteq \mathbb{R} \;\exists L \in \mathbb{R}
\;\exists K \in \mathbb{R}\\
&\quad\Bigl(
    \bigl(x_n \to L \;\land\; x_n \to K\bigr)
    \;\land\;
    L \neq K
\Bigr).
\end{aligned}
\]
\[
\operatorname{ConvergesTo}(x_n,L,\mathbb{R})
\land
\operatorname{ConvergesTo}(x_n,K,\mathbb{R})
\land L\neq K.
\]
\end{description}
\end{remark*}

\begin{remark*}[Interpretation]
The theorem licenses the notation $\displaystyle\lim_{n \to \infty} x_n
= L$: the definite article is justified precisely because the limit,
when it exists, is unique. Without uniqueness the limit would be a
relation rather than a function, and the notation would be ambiguous.
The proof strategy is a separation argument: if $L \neq K$ then
$|L - K| > 0$, and choosing $\varepsilon = |L - K|/2$ produces two
disjoint neighborhoods that the tail of $(x_n)$ cannot simultaneously
occupy. The Hausdorff remark in the failure mode block is the forward
pointer: when sequences are studied in general topological spaces,
uniqueness of limits is equivalent to the Hausdorff separation axiom.
\end{remark*}

\begin{dependencies}
\begin{itemize}
  \item \hyperref[def:convergent-sequence]{Convergence of a Sequence}
  \item \hyperref[thm:triangle-inequality]{Triangle inequality}
\end{itemize}
\end{dependencies}


\begin{theorembox}{Theorem (Limit Preserves Eventual Order)}
% ---------------------------------------------------------
% Theorem: Limit Preserves Eventual Order
% Label: thm:limit-preserves-eventual-order
% ---------------------------------------------------------
\begin{theorem}[Limit Preserves Eventual Order]
\label{thm:limit-preserves-eventual-order}

Let $(x_n)$ and $(y_n)$ be real sequences, and let $L,M\in\mathbb{R}$.
Suppose that $x_n\to L$ and $y_n\to M$. If there exists
$N_0\in\mathbb{N}$ such that $x_n\le y_n$ for every $n\ge N_0$, then
$L\le M$.

\smallskip
\noindent
\hyperref[prf:limit-preserves-eventual-order]{\textit{Go to proof.}}
\end{theorem}
\end{theorembox}

\begin{remark*}[Standard quantified statement]
\[
\begin{aligned}
&\forall (x_n),(y_n)\subseteq\mathbb{R}\;\forall L,M\in\mathbb{R}\\
&\quad\Bigl(
  x_n\to L
  \;\land\;
  y_n\to M
  \;\land\;
  \exists N_0\in\mathbb{N}\;\forall n\ge N_0\;(x_n\le y_n)
\Bigr)\\
&\quad\Longrightarrow\quad L\le M.
\end{aligned}
\]
\end{remark*}

\begin{remark*}[Predicate reading]
\[
\bigl(
\operatorname{ConvergesTo}(x_n,L,\mathbb{R})
\land
\operatorname{ConvergesTo}(y_n,M,\mathbb{R})
\land
\operatorname{Eventually}(x_n,\;x_n\le y_n)
\bigr)
\Longrightarrow
L\le M.
\]
\end{remark*}

\begin{remark*}[Interpretation]
An inequality that holds on a tail survives passage to the limit, but only as
a weak inequality. Strict inequalities on the terms may collapse at the
limit, as in $0<x_n$ with $x_n\to 0$.
\end{remark*}

\begin{dependencies}
\begin{itemize}
  \item \hyperref[def:convergent-sequence]{Convergence of a Sequence}
  \item \hyperref[thm:uniqueness-of-limits]{Uniqueness of Limits}
\end{itemize}
\end{dependencies}


\begin{theorembox}{Theorem (Strict Limit Separation Gives Eventual Order)}
% ---------------------------------------------------------
% Theorem: Strict Limit Separation Gives Eventual Order
% Label: thm:strict-limit-separation-gives-eventual-order
% ---------------------------------------------------------
\begin{theorem}[Strict Limit Separation Gives Eventual Order]
\label{thm:strict-limit-separation-gives-eventual-order}

Let $(x_n)$ and $(y_n)$ be real sequences, and let $A,B\in\mathbb{R}$.
Suppose that $x_n\to A$ and $y_n\to B$. If $A<B$, then there exists
$N\in\mathbb{N}$ such that $x_n<y_n$ for every $n\ge N$.

\smallskip
\noindent
\hyperref[prf:strict-limit-separation-gives-eventual-order]{\textit{Go to proof.}}
\end{theorem}
\end{theorembox}

\begin{remark*}[Standard quantified statement]
\[
\begin{aligned}
&\forall (x_n),(y_n)\subseteq\mathbb{R}\;\forall A,B\in\mathbb{R}\\
&\quad\Bigl(
  x_n\to A
  \;\land\;
  y_n\to B
  \;\land\;
  A<B
\Bigr)\\
&\quad\Longrightarrow\quad
\exists N\in\mathbb{N}\;\forall n\ge N\;(x_n<y_n).
\end{aligned}
\]
\end{remark*}

\begin{remark*}[Predicate reading]
\[
\bigl(
\operatorname{ConvergesTo}(x_n,A,\mathbb{R})
\land
\operatorname{ConvergesTo}(y_n,B,\mathbb{R})
\land
A<B
\bigr)
\Longrightarrow
\operatorname{Eventually}(x_n,\;x_n<y_n).
\]
\end{remark*}

\begin{remark*}[Interpretation]
If the limiting values have a positive gap, then sufficiently late terms
inherit the same order. The proof uses the open interval between $A$ and $B$:
eventually $(x_n)$ lies near $A$ and $(y_n)$ lies near $B$, with room left
between the two tails.
\end{remark*}

\begin{dependencies}
\begin{itemize}
  \item \hyperref[def:convergent-sequence]{Convergence of a Sequence}
  \item \hyperref[thm:uniqueness-of-limits]{Uniqueness of Limits}
\end{itemize}
\end{dependencies}


\begin{theorembox}{Theorem (Eventual Strict Comparison Preserves Weak Limit Order)}
% ---------------------------------------------------------
% Theorem: Eventual Strict Comparison Preserves Weak Limit Order
% Label: thm:eventual-strict-comparison-preserves-weak-limit-order
% ---------------------------------------------------------
\begin{theorem}[Eventual Strict Comparison Preserves Weak Limit Order]
\label{thm:eventual-strict-comparison-preserves-weak-limit-order}

Let $(x_n)$ and $(y_n)$ be real sequences, and let $A,B\in\mathbb{R}$.
Suppose that $x_n\to A$ and $y_n\to B$.
\begin{enumerate}[label=(\alph*)]
  \item If there exists $N\in\mathbb{N}$ such that $x_n<y_n$ for every
  $n\ge N$, then $A\le B$.
  \item If there exists $N\in\mathbb{N}$ such that $x_n>y_n$ for every
  $n\ge N$, then $A\ge B$.
\end{enumerate}

\smallskip
\noindent
\hyperref[prf:eventual-strict-comparison-preserves-weak-limit-order]{\textit{Go to proof.}}
\end{theorem}
\end{theorembox}

\begin{remark*}[Standard quantified statement]
\[
\begin{aligned}
&\forall (x_n),(y_n)\subseteq\mathbb{R}\;\forall A,B\in\mathbb{R}\\
&\quad\Bigl(
  x_n\to A
  \;\land\;
  y_n\to B
  \;\land\;
  \exists N\in\mathbb{N}\;\forall n\ge N\;(x_n<y_n)
\Bigr)
\Longrightarrow A\le B,\\
&\quad\Bigl(
  x_n\to A
  \;\land\;
  y_n\to B
  \;\land\;
  \exists N\in\mathbb{N}\;\forall n\ge N\;(x_n>y_n)
\Bigr)
\Longrightarrow A\ge B.
\end{aligned}
\]
\end{remark*}

\begin{remark*}[Predicate reading]
\[
\begin{aligned}
&\bigl(
\operatorname{ConvergesTo}(x_n,A,\mathbb{R})
\land
\operatorname{ConvergesTo}(y_n,B,\mathbb{R})
\land
\operatorname{Eventually}(x_n,\;x_n<y_n)
\bigr)
\Longrightarrow A\le B,\\
&\bigl(
\operatorname{ConvergesTo}(x_n,A,\mathbb{R})
\land
\operatorname{ConvergesTo}(y_n,B,\mathbb{R})
\land
\operatorname{Eventually}(x_n,\;x_n>y_n)
\bigr)
\Longrightarrow A\ge B.
\end{aligned}
\]
\end{remark*}

\begin{remark*}[Interpretation]
Strict comparison of terms does not force strict comparison of limits. The
strict inequality can collapse in the limit, so the conclusion is weak:
for example $1/n>0$ for every $n\in\mathbb{N}$, but both sides have limit
$0$.
\end{remark*}

\begin{dependencies}
\begin{itemize}
  \item \hyperref[thm:limit-preserves-eventual-order]{Limit Preserves Eventual Order}
\end{itemize}
\end{dependencies}


\begin{theorembox}{Theorem (Constant Comparison for Sequence Limits)}
% ---------------------------------------------------------
% Theorem: Constant Comparison for Sequence Limits
% Label: thm:constant-comparison-sequence-limits
% ---------------------------------------------------------
\begin{theorem}[Constant Comparison for Sequence Limits]
\label{thm:constant-comparison-sequence-limits}

Let $(x_n)$ be a real sequence, and let $A,B\in\mathbb{R}$. Suppose
that $x_n\to A$.
\begin{enumerate}[label=(\alph*)]
  \item If there exists $N\in\mathbb{N}$ such that $x_n\le B$ for every
  $n\ge N$, then $A\le B$.
  \item If there exists $N\in\mathbb{N}$ such that $x_n<B$ for every
  $n\ge N$, then $A\le B$.
  \item If there exists $N\in\mathbb{N}$ such that $x_n\ge B$ for every
  $n\ge N$, then $A\ge B$.
  \item If there exists $N\in\mathbb{N}$ such that $x_n>B$ for every
  $n\ge N$, then $A\ge B$.
\end{enumerate}

\smallskip
\noindent
\hyperref[prf:constant-comparison-sequence-limits]{\textit{Go to proof.}}
\end{theorem}
\end{theorembox}

\begin{remark*}[Standard quantified statement]
\[
\begin{aligned}
&\forall (x_n)\subseteq\mathbb{R}\;\forall A,B\in\mathbb{R}\\
&\quad\Bigl(
  x_n\to A
  \;\land\;
  \exists N\in\mathbb{N}\;\forall n\ge N\;(x_n\le B)
\Bigr)
\Longrightarrow A\le B,\\
&\quad\Bigl(
  x_n\to A
  \;\land\;
  \exists N\in\mathbb{N}\;\forall n\ge N\;(x_n<B)
\Bigr)
\Longrightarrow A\le B,\\
&\quad\Bigl(
  x_n\to A
  \;\land\;
  \exists N\in\mathbb{N}\;\forall n\ge N\;(x_n\ge B)
\Bigr)
\Longrightarrow A\ge B,\\
&\quad\Bigl(
  x_n\to A
  \;\land\;
  \exists N\in\mathbb{N}\;\forall n\ge N\;(x_n>B)
\Bigr)
\Longrightarrow A\ge B.
\end{aligned}
\]
\end{remark*}

\begin{remark*}[Predicate reading]
\[
\begin{aligned}
&\operatorname{ConvergesTo}(x_n,A,\mathbb{R})
\land
\operatorname{Eventually}(x_n,\;x_n\le B)
\Longrightarrow A\le B,\\
&\operatorname{ConvergesTo}(x_n,A,\mathbb{R})
\land
\operatorname{Eventually}(x_n,\;x_n<B)
\Longrightarrow A\le B,\\
&\operatorname{ConvergesTo}(x_n,A,\mathbb{R})
\land
\operatorname{Eventually}(x_n,\;x_n\ge B)
\Longrightarrow A\ge B,\\
&\operatorname{ConvergesTo}(x_n,A,\mathbb{R})
\land
\operatorname{Eventually}(x_n,\;x_n>B)
\Longrightarrow A\ge B.
\end{aligned}
\]
\end{remark*}

\begin{remark*}[Interpretation]
This is the constant version of the comparison theorem. A fixed real number
$B$ may be treated as the constant sequence with value $B$, so each conclusion
is a direct application of eventual order preservation. Again, strict term
inequalities produce only weak inequalities between limits.
\end{remark*}

\begin{dependencies}
\begin{itemize}
  \item \hyperref[def:constant-sequence]{Constant Sequence}
  \item \hyperref[thm:limit-preserves-eventual-order]{Limit Preserves Eventual Order}
  \item \hyperref[thm:eventual-strict-comparison-preserves-weak-limit-order]{Eventual Strict Comparison Preserves Weak Limit Order}
\end{itemize}
\end{dependencies}


\begin{theorembox}{Theorem (Constant Squeeze Theorem)}
% ---------------------------------------------------------
% Theorem: Constant Squeeze Theorem
% Label: thm:constant-squeeze-theorem
% ---------------------------------------------------------
\begin{theorem}[Constant Squeeze Theorem]
\label{thm:constant-squeeze-theorem}

Let $(x_n)$ be a real sequence, and let $L\in\mathbb{R}$. If there
exists $N_0\in\mathbb{N}$ such that $L\le x_n\le L$ for every
$n\ge N_0$, then $x_n\to L$.

\smallskip
\noindent
\hyperref[prf:constant-squeeze-theorem]{\textit{Go to proof.}}
\end{theorem}
\end{theorembox}

\begin{remark*}[Standard quantified statement]
\[
\begin{aligned}
&\forall (x_n)\subseteq\mathbb{R}\;\forall L\in\mathbb{R}\\
&\quad\Bigl(
  \exists N_0\in\mathbb{N}\;\forall n\ge N_0\;(L\le x_n\le L)
\Bigr)\\
&\quad\Longrightarrow\quad
\forall \varepsilon>0\;\exists N\in\mathbb{N}\;\forall n\ge N\;
\bigl(|x_n-L|<\varepsilon\bigr).
\end{aligned}
\]
\end{remark*}

\begin{remark*}[Predicate reading]
\[
\operatorname{Eventually}(x_n,\;L\le x_n\le L)
\Longrightarrow
\operatorname{ConvergesTo}(x_n,L,\mathbb{R}).
\]
\end{remark*}

\begin{remark*}[Interpretation]
This is the degenerate squeeze theorem: the lower and upper barriers are the
same constant sequence. It says that eventual equality to a real number is
enough for convergence to that real number.
\end{remark*}

\begin{dependencies}
\begin{itemize}
  \item \hyperref[def:convergent-sequence]{Convergence of a Sequence}
  \item \hyperref[def:constant-sequence]{Constant Sequence}
\end{itemize}
\end{dependencies}


\begin{theorembox}{Theorem (Sequence Squeeze Theorem)}
% ---------------------------------------------------------
% Theorem: Sequence Squeeze Theorem
% Label: thm:sequence-squeeze-theorem
% ---------------------------------------------------------
\begin{theorem}[Sequence Squeeze Theorem]
\label{thm:sequence-squeeze-theorem}

Let $(a_n)$, $(x_n)$, and $(b_n)$ be real sequences, and let
$L\in\mathbb{R}$. Suppose that $a_n\to L$ and $b_n\to L$. If there
exists $N_0\in\mathbb{N}$ such that $a_n\le x_n\le b_n$ for every
$n\ge N_0$, then $x_n\to L$.

\smallskip
\noindent
\hyperref[prf:sequence-squeeze-theorem]{\textit{Go to proof.}}
\end{theorem}
\end{theorembox}

\begin{remark*}[Standard quantified statement]
\[
\begin{aligned}
&\forall (a_n),(x_n),(b_n)\subseteq\mathbb{R}\;\forall L\in\mathbb{R}\\
&\quad\Bigl(
  a_n\to L
  \;\land\;
  b_n\to L
  \;\land\;
  \exists N_0\in\mathbb{N}\;\forall n\ge N_0\;(a_n\le x_n\le b_n)
\Bigr)\\
&\quad\Longrightarrow\quad
\forall \varepsilon>0\;\exists N\in\mathbb{N}\;\forall n\ge N\;
\bigl(|x_n-L|<\varepsilon\bigr).
\end{aligned}
\]
\end{remark*}

\begin{remark*}[Predicate reading]
\[
\bigl(
\operatorname{ConvergesTo}(a_n,L,\mathbb{R})
\land
\operatorname{ConvergesTo}(b_n,L,\mathbb{R})
\land
\operatorname{Eventually}(x_n,\;a_n\le x_n\le b_n)
\bigr)
\Longrightarrow
\operatorname{ConvergesTo}(x_n,L,\mathbb{R}).
\]
\end{remark*}

\begin{remark*}[Interpretation]
The middle sequence is trapped between two tails that approach the same
number. Once both barriers lie inside an $\varepsilon$-neighborhood of $L$,
the trapped sequence must also lie inside that neighborhood.
\end{remark*}

\begin{dependencies}
\begin{itemize}
  \item \hyperref[def:convergent-sequence]{Convergence of a Sequence}
  \item \hyperref[thm:limit-preserves-eventual-order]{Limit Preserves Eventual Order}
\end{itemize}
\end{dependencies}


\begin{theorembox}{Theorem (Absolute-Value Squeeze Theorem)}
% ---------------------------------------------------------
% Theorem: Absolute-Value Squeeze Theorem
% Label: thm:absolute-value-squeeze-theorem
% ---------------------------------------------------------
\begin{theorem}[Absolute-Value Squeeze Theorem]
\label{thm:absolute-value-squeeze-theorem}

Let $(x_n)$ and $(u_n)$ be real sequences, and let $L\in\mathbb{R}$.
Suppose that $u_n\to 0$. If there exists $N_0\in\mathbb{N}$ such that
$|x_n-L|\le u_n$ for every $n\ge N_0$, then $x_n\to L$.

\smallskip
\noindent
\hyperref[prf:absolute-value-squeeze-theorem]{\textit{Go to proof.}}
\end{theorem}
\end{theorembox}

\begin{remark*}[Standard quantified statement]
\[
\begin{aligned}
&\forall (x_n),(u_n)\subseteq\mathbb{R}\;\forall L\in\mathbb{R}\\
&\quad\Bigl(
  u_n\to 0
  \;\land\;
  \exists N_0\in\mathbb{N}\;\forall n\ge N_0\;(|x_n-L|\le u_n)
\Bigr)\\
&\quad\Longrightarrow\quad
\forall \varepsilon>0\;\exists N\in\mathbb{N}\;\forall n\ge N\;
\bigl(|x_n-L|<\varepsilon\bigr).
\end{aligned}
\]
\end{remark*}

\begin{remark*}[Predicate reading]
\[
\bigl(
\operatorname{ConvergesTo}(u_n,0,\mathbb{R})
\land
\operatorname{Eventually}(x_n,\;|x_n-L|\le u_n)
\bigr)
\Longrightarrow
\operatorname{ConvergesTo}(x_n,L,\mathbb{R}).
\]
\end{remark*}

\begin{remark*}[Interpretation]
This is the estimate form of the squeeze theorem. Instead of constructing
two explicit bounding sequences, it bounds the error $|x_n-L|$ by a null
sequence.
\end{remark*}

\begin{dependencies}
\begin{itemize}
  \item \hyperref[def:convergent-sequence]{Convergence of a Sequence}
  \item \hyperref[def:null-sequence]{Null Sequence}
  \item \hyperref[thm:sequence-squeeze-theorem]{Sequence Squeeze Theorem}
\end{itemize}
\end{dependencies}

